\section{Task 5: Full Problem with Quality Constraints}
\label{sec:task5}

\subsection{Problem Setup}
\label{sec:task5-setup}
This task uses the same production volume and colour as Task 3, but adds minimum average quality constraints. The setup is 5000 L amber beer with
\[
Q^S_{\min}=3,\qquad Q^H_{\min}=3,\qquad Q^Y_{\min}=3.
\]

\begin{table}[H]
\renewcommand{\arraystretch}{1.15}
\caption{Task 5 Problem Setup}
\label{tab:task35-setup}
\centering
\small
\begin{tabular}{@{}ll@{}}
\toprule
Parameter & Value \\
\midrule
Beer volume & 5000 L \\
Colour & amber \\
EBC bounds & 20 to 29 \\
Malt-source quality & at least 3 \\
Hop quality & at least 3 \\
Yeast quality & at least 3 \\
\bottomrule
\end{tabular}
\end{table}

\subsection{Optimal Solution}
\label{sec:task5-solution}
The optimal quality-constrained solution uses process choice $(z_{\mathrm{malt}},z_{\mathrm{mash}})=(0,1)$. The total cost is 4891.923 and the achieved EBC is 20.000. All three achieved group qualities are 3.000.

\begin{table}[H]
\renewcommand{\arraystretch}{1.15}
\caption{Task 5 Optimal Solution Summary}
\label{tab:task35-solution}
\centering
\small
\begin{tabular}{@{}ll@{}}
\toprule
Item & Result \\
\midrule
Total cost & 4891.923 \\
Process choice & $(0,1)$ \\
Achieved EBC & 20.000 \\
Malt-source quality & 3.000 \\
Hop quality & 3.000 \\
Yeast quality & 3.000 \\
\bottomrule
\end{tabular}
\end{table}

\begin{table}[H]
\renewcommand{\arraystretch}{1.15}
\caption{Task 5 Non-Zero Ingredients}
\label{tab:task35-ingredients}
\centering
\small
\begin{tabular}{@{}llr@{}}
\toprule
Group & Variable & Amount (kg) \\
\midrule
Malt & \texttt{malt03} & 901.457 \\
Malt & \texttt{malt06} & 17.497 \\
Malt & \texttt{malt12} & 34.994 \\
Hops & \texttt{hop07} & 3.421 \\
Hops & \texttt{hop10} & 3.421 \\
Yeast & \texttt{yeast03} & 3.750 \\
\bottomrule
\end{tabular}
\end{table}

\subsection{Comparison with Task 3}
\label{sec:task5-comparison}
The quality constraints increase the cost from 4497.687 in Task 3 to 4891.923 in Task 5. This is expected because the feasible region becomes smaller. The model can no longer use only the cheapest low-quality combination from Task 3; it must choose ingredients whose weighted averages satisfy the three quality lower bounds. The optimal EBC also moves from the upper amber boundary in Task 3 to the lower amber boundary in Task 5.

% Values in this section come from task_35.ipynb.
